\chapter*{General Conclusion}
\addcontentsline{toc}{chapter}{General Conclusion}

This internship transformed a concrete production-information delay into a complete Industrial Engineering prototype. Soluble MAP moisture is confirmed by laboratory analysis approximately every two hours while dryer conditions evolve continuously. The response was not to replace the laboratory or control system, but to organize an intermediate decision-support layer: residence-aware moisture prediction, row-level novelty detection, subsystem-oriented diagnosis, traceable PostgreSQL storage and a centralized Power BI supervision interface.

The final implementation is technically coherent. The supplied deterministic 92-day source contains 1,589,760 process rows and 1,104 sparse laboratory observations. It feeds an executable Notebook 01--04 chain, prevents future information from entering the aligned features, and yields 1,103 rows with the original 16 named moisture-model inputs. Direct process values use the residence-adjusted time; density and final product temperature use the most recent strictly previous laboratory result. Notebook 03 selected Ridge using a 772/166/165 chronological TRAIN/VALIDATION/TEST split; TEST MAE is 0.001069, RMSE 0.001403 and $R^2=0.824530$. Notebook 04 exports the 15-feature anomaly model, scaler and reference profile. The active service loads those exact artifacts while replaying the held-out TEST interval into PostgreSQL and the Power BI views. The five-second interval is a selected prototype replay/visualization interval, not native plant acquisition. The supplied-file hash, artifact hashes, automated tests and live SQL/TMDL validation make these boundaries machine-checkable.

The decisive conclusion remains the scope boundary. The demonstration data are synthetic, the service is not connected to live PCS7, statistical limits are not approved plant limits, and no autonomous action is implemented. Production value can only be assessed after retraining and shadow validation on representative OCP historical data, with operator involvement, industrial cybersecurity, alarm governance and robust operations engineering.

The final classification is therefore exact: this deliverable is an \textbf{engineering prototype, not plant-ready}.

From an Operational Excellence perspective, the project shows how data and AI tools can strengthen quality visibility, supervision discipline and problem-solving speed when they are integrated with process understanding and explicit human authority. Its strongest deliverable is therefore not an isolated model, but a transparent architecture and validation path that can support a future OCP industrial study.
